% Spellchecking done with grammarly
\section{Baseline Methodology}
The main starting point for this thesis is the papers Visual Fingerprints of Vibration Signals
Using Time Delay Embedding. The goal of this paper was to create an algorithm which solves the issue of plotting vibration data on line graph where you can't distinguish between noise and the actual vibrations. The overall goal was to: further enhance and complement the visual inspection in the time
domain, enabling users to quickly gain an overview of different
vibration profiles within the signal or to detect periodicity within
noise. (\cite{rakuschek2025visual} p.1) In the background this was done by using  Time Delay Embeddings
(TDE), which is an way of taking the time series data and transform it into a point cloud that can be projected onto a 2D plane. Another really important aspect of the algorithm which is to be transformed is the use of Principal Component Analysis which is a dimensionality reduction technique and helps to reduce the number of features in a dataset while keeping the most important information. It changes complex datasets by transforming correlated features into a smaller set of uncorrelated components (\cite{geeksforgeeks2025pca}). When this two algorithms are applied on a harmonic oscillation the result is a ellipsoid. When noise is added its really hard to distinguished between noise or the hidden signal. The vibration signal is then revealed by the TDE with different window sizes. This behavior is visualized by the appendix 1 which shows the contrast between time series data and the suggested approach which resembles a human fingerprint.
When noise is introduced, distinguishing the underlying hidden signal from the noise becomes increasingly difficult. However, applying Time Delay Embedding (TDE) with appropriate window sizes effectively reveals the underlying vibration signal.